\chapter{Mynotes from 2026 class}
Some references used here:
\begin{itemize}
    \item Mumford's red book~\cite{mumford1988red}
    \item Humphreys Linear algebraic groups~\cite{humphreys1975linear}
\end{itemize}

\section{Lecture 1 and 2}
First some definitions and lemmas worth to remember.
\begin{itemize}
    \item (Integral extension). Let $R$ be a ring and $S\subset R$ a subring. We say that 
    $R$ is integrally dependent over $S$ if every $r\in R$ is integral over $S$, that is, there 
    exist a monic polynomial $f\in S[X]$ such that $f(r)=0$.
    \begin{lemma}\label{lemma : characterization of integral extensions}
        Let $S\subset R$ be rings. 
        A finite subset of elements $b_1,\dots, b_m\in R$ are all integral over $S$ if and only if 
        $S[b_1,\dots, b_m]$ is finitely generated (as a $S-$module) over $S$.
    \end{lemma}
    Proof: $(\Rightarrow)$ Let $b\in R$ be integral over $S$, i.e., there exist 
    \begin{equation}
     f =    X^d + a_1X^{d-1} +\cdots + a_d \in S[X]
    \end{equation}
    such that $f(b)=0$. Then, using long division, we see that $S[b]=S +\cdots + S b^{d-1}$ which is finitely generated. 
    Now, if $b_1,\dots, b_m$ are all integral over $S$, we can argue by induction that $S[b_1,\dots, b_{m-1}]$ is finitely generated 
    over $S$. However, since $b_m$ is integral over $S[b_1,\dots, b_{m-1}],\;S[b_1,\dots,b_m]$ is finitely generated 
    over $[b_1,\dots, b_{m-1}]$ and it must be finitely generated over $S$ aswell.\par  
    $(\Leftarrow)$ Suppose that $S[b_1,\dots, b_m] = S\beta_1+\cdots+S\beta_n$ for some $\beta_i\in R$. Then for any 
    $b\in R$ which fixes $S[b_1,\dots,b_m]$, we have $b\beta_i= \sum_j a_{ij}\beta_j$ for some $(a_{ij})\in\text{Mat}_n(S)$. 
    Therefore 
    \begin{equation}
      0 =  (bI-(a_{ij}))^*(bI-(a_{ij}))v = \det(bI-(a_{ij}))v
    \end{equation}
    where $v=(\beta_1,\dots,\beta_n)^T$. Therefore $\det(bI-(a_{ij}))\beta_i=0,\forall i$ and since 
    $1\in S\beta_1+\cdots+S\beta_n$, we see that $\det(bI-(a_{ij}))=0$. Since $\det(X-(a_{ij}))\in S[X]$ is monic, we 
    are done.\qed 
    \begin{corollary}\label{corollary : transitive property integral dependence}
        (Transitive property of integral dependence). Given rings $R \supset S\supset T$, if $R$ is integral over $S$ and 
        $S$ is integral over $T$, then $R$ is integral over $T$.\par 
        \textcolor{red}{Question: Do we need Noetherian hypothesis using this statement?}
    \end{corollary}
    Proof: Let $r\in R$, then by hypothesis there exist $f=X^d+s_1X^{d-1}+\cdots\in S[X]$ such that $f(r)=0$.
    Define $A=T[s_1,\dots,s_d]$, then $A[r]$ is a finitely generated $A-$module because $S\subset A$ and lemma~(\ref{lemma : characterization of integral extensions}). 
    Now, since $A$ is a finitely generated $T-$module by lemma~(\ref{lemma : characterization of integral extensions}),
    we see that $A[r]$ is a finitely generated $T-$module. Now, \textcolor{red}{assuming $T$ is Noetherian} we see that 
    $T[r]$ is finitely generated as a $T-$module.\qed

    \begin{remark}
        We can certainly say (using the previous proof) that if $r\in R$, then $S[r]$ is finitely generated 
        as a $T-$module.
    \end{remark}

    \begin{lemma}\label{lemma : extension field transcendence degree}
        Let $K\supset k$ ($k$ a perfect field) be a finitely generated extension field with transcendence degree $\partial_k(K)=d$, then 
        there exist generators $z_1,\dots,z_d,z_{d+1}$ such that 
        \begin{itemize}
            \item[(1)] $K=k(z_1,\dots,z_{d+1})$
            \item[(2)] $\{z_1,\dots,z_d\}$ are algebraicaly independent
            \item[(3)] $z_{d+1}$ is separable over $k(z_1,\dots, z_d)$ 
        \end{itemize} 
    \end{lemma}  
\end{itemize}

\begin{dang}
\begin{lemma}
    (Noether's normalization lemma - Mumford's red book). Let $R$ be a finitely generated algebra over $k$. Assume also that $R$ is an integral domain.
    If $R$ has transcende degree $n$ over $k$ (which means that $\partial_k\text{Frac}(R)=n$), then 
    there exist $x_1,\dots, x_n\in R$ algebraicaly independent and such that $R$ is integral 
    over $k[x_1,\dots,x_n]$.
\end{lemma}
\end{dang}
\TODO Study prove this version of the lemma using the statement made in class.

\begin{dang}
    \begin{lemma}\label{lemma : Noether normalization}
        (Noether's normalization lemma - Statement made in class). Let $R$ be a finitely generated 
        $k-$algebra. Then there exist $x_1,\dots,x_n\in R$ algebraically independent over $k$ such that 
        $R$ is integral over $k[x_1,\dots,x_n]$.
    \end{lemma}
\end{dang}
Proof (Nagata): Since $R$ is finitely generated as a $k$-algebra, there exist generators $y_1,\dots,y_m$ such that 
$R= k[y_1,\dots, y_m]$. Now suppose that $y_1,\dots, y_m$ are algebraically independent, then we can take $x_i = y_i$. 
If not, then there exist $p\in k[Y_1,\dots, Y_m]\setminus0$ such that $p(y_1,\dots,y_m)=0$. Write 
\begin{equation}
    p(Y_1,\dots,Y_m)=\sum_{\alpha\in\bbZ^m_\geq0}a_\alpha Y_1^{\alpha_1}\cdots Y_m^{\alpha_m}
\end{equation}
and define a new polynomial 
\begin{equation}
    \begin{split}
        q(Y_1,\dots,Y_m) &= p(Y_1,Y_2+Y_1^{r_2},\dots, Y_m+Y_1^{r_m})\\
        &=\sum_\alpha a_\alpha(Y_1^{\alpha_1+r_2\alpha_2+\cdots+r_m\alpha_m}+\cdots) 
    \end{split}
\end{equation}
for some choice of positive integers $r_2,\dots,r_m$. Take a positive number $b>\alpha_i,\;\forall \alpha_i$ such that 
$a_\alpha\neq0$. If we take $r_i = b^i$, then $\alpha_1+\alpha_2b+\cdots +\alpha_mb^{m-1}$ are all different numbers in base $b$. 
So, from different powers of the form $\alpha_1+r_2\alpha_2+\cdots+r_m\alpha_m$ we have a maximal one 
$d=\max\{\alpha_1+\cdots r_m\alpha_m:a_\alpha\neq0\}$ and 
\begin{equation}
    q(Y_1,\dots,Y_m) = aY_1^d + \text{smaller degree terms} \in k[Y_2,\dots, Y_m][Y_1],\;a\neq0.
\end{equation}
Now, if we let $z_i=y_i - y_1^{r_i}\in R,i=2,\dots,m$, we see that $q(y_1,z_2,\dots,z_m)=p(y_1,\dots,y_m)=0$ and 
$y_1$ is integral over $k[z_2,\dots, z_m]$. By induction on number of generators, there exist $x_1,\dots, x_n$
algebraically independent such that $k[z_2,\dots, z_m]$ is integral over $k[x_1,\dots, x_n]$. So, by lemma~\ref{lemma : characterization of integral extensions}
and its corollary~(\ref{lemma : characterization of integral extensions}), we see that $k[y_1,z_2,\dots,z_m]$ is integral over $k[x_1,\dots,x_m]$. Now, by definition, it's clear 
that $k[y_1,z_2,\dots,z_m]$ is a finitely generated $R-$module. So by the corollary~(\ref{corollary : transitive property integral dependence}) again, we 
see that $R$ is integral over $k[x_1,\dots,x_m]$.\qed 

\begin{dang}
    \begin{lemma}\label{lemma : R is a field, S is a field}
        Let $S\subset R$ be a ring extension. If $R$ is a field, which integral over $S$, then $S$ is a field. 
    \end{lemma}
\end{dang}
Proof: Let $s\in S\setminus0$ and $r\in R$ be its inverse. Since $R$ is integral over $S$, there exist
$f=X^d+a_1X^{d-1}+\cdots\in S[X]$ be such that $f(r)=0$. Therefore 
\begin{equation}
   0= s^{d-1}f(r) = r + a_1 + sa_2+\cdots + s^{d-1}a_0\Rightarrow r\in S.
\end{equation} 
\qed

\begin{remark}
    Let $L|K$ be a field extension. Then
    \begin{center}
        \begin{tikzcd}
            L|K \text{ finite extension}\arrow[r, bend left= 30] &  L|K \text{ algebraic extension}
            \arrow[l, bend left= 30, "\text{+ finitely generated extension}"]
        \end{tikzcd}
    \end{center}
\end{remark}

\begin{dang}
    \begin{theorem}\label{theorem : Zariski 1}
        (Zariski). Let $k$ be a field and $\fr{m}\subset R=k[X_1,\dots, X_m]$ be a maximal ideal. Then $L=R/\fr{m}$ is 
        a finite extension field of $k$.
    \end{theorem}
\end{dang}
Proof: By Noether normalization lemma~(\ref{lemma : Noether normalization}), there exist 
$a_1,\dots, a_n\in L$ algebraically independent such that $L$ is integral over $k[a_1,\dots,a_n]$. By lemma~(\ref{lemma : R is a field, S is a field}), 
we conclude that $k[a_1,\dots,a_n]$ is a field. Suppose that $n>0$, so $a_1$ has an inverse. This means that 
$a_1g(a_1,\dots, a_n) = 1$ for some $g\in k[X_1,\dots, X_n]$, which constradicts the fact that $a_1,\dots, a_n$ are 
alegrbaically independent. Then we conclude that $L$ is an integral (and then algebraic) over $k$. However, since 
$L$ is a finitely generated extension ($L=k(\overline{X}_1,\dots,\overline{X}_m)$) over $k$ we conclude that $L$ is in fact 
a finite field extension over $k$.\qed

\begin{dang}
    \begin{theorem}\label{theorem : Hilbert's zero theorem}
        Suppose that $k$ is an algebraically closed field.
        Let $I\subset R =k[X_1,\dots, X_m]$ be an ideal, then $V(I)=\varnothing\iff I = k[X_1,\dots, X_m]$.
    \end{theorem}
\end{dang}
Proof: If $I\neq R$, then there exist a maximal ideal $\fr{m}\supset I$. So $V(I)\supset V(\fr{m})$ and it's sufficient to 
prove that $V(\fr{m})\neq\varnothing$. By theorem~(\ref{theorem : Zariski 1}), we have that $L=R/\fr{m}$ is a finite 
$k-$extension, so we must have $L=k$ because $k$ is algebraically closed. So there exist $a_i\in k,i=1,\dots,m$ such that 
$X_i - a_i\in\fr{m}$ and $(X_1-a_1,\dots, X_m-a_m)\subset \fr{m}$. However $\varphi: R\to k:F\mapsto F(a_1,\dots,a_m)$ has kernel 
$\Ker(\varphi)=(X_1-a_1,\dots, X_m-a_m)$ and since $k$ is a field, $R/\Ker(\varphi)$ must be  a field and 
$\Ker(\varphi)$ is a maximal ideal. Therefore we conclude that $\fr{m}=(X_1-a_1,\dots, X_m-a_m)$. Finally
$V(X_1-a_1,\dots,X_m-a_m)=\{a_1,\dots,a_m\}\neq\varnothing$.\qed

\begin{remark}
    The hypothesis $k=\overline{k}$ in the previous theorem is crucial. For example if $k=\bbR$, we have 
    $V(x^2+1)=\varnothing$.
\end{remark}

\begin{corollary}\label{corollary : maximal ideals in k[x1,...,xn] ring characterization}
    From the proof of the previous theorem, we see that there exist a bijection 
    \begin{equation}
        \begin{split}
            k^n &\longleftrightarrow \left\{\begin{matrix}\text{maximal ideals}\\\fr{m}\subset R\end{matrix}\right\}\\
           (a_1,\dots,a_n)&\longmapsto (X_1-a_1,\dots, X_n-a_n)
        \end{split}
    \end{equation}
\end{corollary}


\begin{dang}
    \begin{theorem}\label{theorem : Hilbert Nullstelensatz}
        (Hilbert Nullstelensatz). Let $A\subset k[X_1,\dots, X_n]$ be an ideal, then 
        \begin{equation}
            IV(A) = \sqrt{A}.
        \end{equation}
    \end{theorem}
\end{dang}
Proof: The inclusion $\sqrt{A}\subset IV(A)$ is easy. Now, assume $g\in IV(A)$, where $A=(f_1,\dots, f_m)$. This 
means that we have the hypothesis:
\begin{equation}
    f_i(a_1,\dots, a_n)=0,\;\forall i=1,\dots,m\Rightarrow g(a_1,\dots, a_n)=0\; (*)
\end{equation}
Now, define the ideal 
\begin{equation}
    I = Ak[X_1,\dots, X_{n+1}] + (1- g(X_1,\dots, X_n)X_{n+1}).
\end{equation}
If $I$ is a proper ideal, then by corollary~(\ref{corollary : maximal ideals in k[x1,...,xn] ring characterization}), we have 
$I\subset (X_1-a_1,\dots, X_{n+1}-a_{n+1})$, which is the kernel of the evaluation map 
$k[X_1,\dots,X_{n+1}]\rightarrow k: f\mapsto f(a_1,\dots, a_{n+1})$. This means that 
\begin{equation}
    \left\{\begin{aligned}
        &f_i(a_1,\dots,a_n)=0,\;\forall i=1,\dots,m\\
        &g(a_1,\dots,a_n)a_{n+1}=1
    \end{aligned}\right.
\end{equation}
which contradicts our initial hypothesis $(*)$. Therefore, $I$ is not a proper ideal, so $1\in I$ and there exists 
$h_1,\dots, h_{m+1}\in k[X_1,\dots, X_{n+1}]$ such that 
\begin{equation}
    1 = \sum_{j=1}^mh_j(X_1,\dots, X_{n+1})f_j(X_1,\dots,X_n)  + h_{m+1}(X_1,\dots,X_{n+1})(1-g(X_1,\dots,X_n)X_{n+1}).
\end{equation}
Now, by projecting this equation onto $k[X_1,\dots, X_{n+1}]/(1-g X_{n+1})$, we get 
\begin{equation}
    g^\ell(X_1,\dots,X_n) = \sum_{j=1}^m h_j(X_1,\dots, X_{n+1})g^\ell  \times f_j(X_1,\dots,X_n)
\end{equation}
modulo $(1-g X_{n+1})$, where $H_j = h_j(X_1,\dots, X_{n+1})g^\ell\in k[X_1,\dots,X_n],\;\forall j$ if $\ell\gg0$. 
Since $\{\overline{X}_1,\dots,\overline{X}_{n}\}$ is algebraically free in $k[X_1,\dots, X_{n+1}]/(1-g X_{n+1})$, we 
conclude that, in fact, we have equality 
\begin{equation}
     g^\ell(X_1,\dots,X_n) = \sum_{j=1}^m H_j(X_1,\dots,X_n)  \times f_j(X_1,\dots,X_n)
\end{equation}
in $k[X_1,\dots,X_n]$.\qed

\begin{remark}
    We could have used the homomorphism 
    \begin{equation}
        k[X_1,\dots,X_{n+1}]\rightarrow k(X_1,\dots,X_n):X_{n+1}\mapsto g,\;X_i\mapsto X_i,\;i\neq n+1
    \end{equation}
    at the end of the previous proof.
\end{remark}

\section{Lecture 3}
\begin{dang}
    \begin{definition}
        We define an affine algebraic set $V\subset k^n$ to be a set of the form $V=V(S)$ for 
        some $S\subset k[\underline{X}]$.
    \end{definition}
\end{dang}

\begin{notation}
    Let $V$ be an affine algebraic set. We write 
    \begin{equation}
        D_V(g) = \{p\in V: g(p)\neq0\} = V\setminus V(g).
    \end{equation}
    These open sets are called ``principal open subsets''.
\end{notation}

\begin{dang}
    \begin{lemma}
        Any open subset $U\subset V$ of an affine algebraic set is quasicompact.
    \end{lemma}
\end{dang}
Proof: Consider a covering $U=\bigcup_{\lambda\in\Lambda}U_\lambda$.
Observe that $V\setminus U=V(g_1,\dots, g_m)$, so we have the equality $U=\bigcup_{i=1}^m D_V(g_i)$. In particular, we see 
that principal open subsets form a basis for the Zariski topology on $V$. Therefore, 
it's sufficient to prove that a principal open subset is quasicompact. So we assume $U=D_V(g)$. Furthermore, we have 
$U_\lambda = \bigcup_{i=1}^{n_\lambda}D_V(g_{\lambda,i})$ and 
\begin{equation}
    U=\bigcup_{\lambda\in\Lambda}\bigcup_{i=1}^{n_\lambda}D_V(g_{\lambda,i}).
\end{equation}
So, without lost of generalizaiton, we can assume $U=\bigcup_{\lambda\in\Lambda}D_V(g_\lambda)$ from beginning. Now 
\begin{equation}
    \begin{split}
        D_V(g) & = \bigcup_{\lambda\in\Lambda}D_V(g_\lambda)\iff\\
      V\cap V(g) &= V\cap V(\langle g_\lambda:\lambda\in\Lambda\rangle)\iff\\ 
      V\cap V(g) &= V\cap V(\langle g_\lambda:\lambda\in I\rangle)\;(I\subset \Lambda\;\text{ finite}).
    \end{split}
\end{equation}
In the last equation we used Hilbert basis theorem: The sequence
\begin{equation}
    \langle g_{\lambda_1} \rangle \subset  \langle g_{\lambda_1},g_{\lambda_2} \rangle \subset \cdots 
\end{equation}
must terminate at some point, because $k[\underline{X}]$ is Noetherian. Then we conclude that 
$D_V(g)=\bigcup_{\lambda\in I}D_V(g_\lambda)$.\qed


\begin{dang}
    \begin{lemma}
        Suppose that $k$ is an infinity field, then open sets $U\subset k^n$ are dense in the Zariski topology.
    \end{lemma}
\end{dang}
Proof: Since basic open sets form a basis for the Zariski topology on $k^n$, it's sufficient to prove that 
$D(f)$ is dense in $k^n$. Let $\overline{D(f)}=V(g_1,\dots, g_m),\;m\geq1$ and $x\in k^n$, then we have 
for all $i=1,\dots, m$:
\begin{equation}
    f(x)\neq0\Rightarrow g_i(x) = 0\Rightarrow fg_i(x)=0,\;\forall x\in k^n.
\end{equation}
Therefore, since $k$ is infinity, we have $fg_i=0\in k[\underline{X}],\;\forall i$ and $g_i=0,\;\forall i$. Then we see that 
$\overline{D(f)}=k^n$.\qed


\begin{dang}
    \begin{definition}
    (Regular functions). Let $V\subset k^n$ be an affine algebraic set and $U\subset V$ be an open set. 
    A map $\varphi:U\to k$ is said to be \textit{regular} if for any $p\in U$, there exist $p\in U^\prime\subset U$, 
    polynomials $f,g\in k[\underline{X}]$ such that $g(q)\neq0,\;\forall q\in U^\prime$ and  
    $\varphi|_{U^\prime}=f/g|_{U^\prime}$.    
    \end{definition}
\end{dang}

\begin{dang}
    \begin{notation}
    Let $V$ be an affine algebraic set and let $U\subset V$ be an open set, we denote 
    $\mathcal{O}_V(U)=\{U\overset{\text{regular}}{\rightarrow}k\}$.
\end{notation}
\end{dang}


\begin{dang}
    \begin{lemma}\label{lemma: regular functions on principal open sets}
        Let $V\subset k^n$ be an affine algebraic set and take $g\in k[X_1,\dots, X_n]\setminus I(Y)$, 
        so $D_V(g)\neq\varnothing$. Then the set of regular functions 
        $\varphi: D_V(g)\to k$ is of the form $f/g^m|_{D_V(g)}$ for some $f\in k[X_1,\dots, X_n]$ and $m>0$. Then we have  
        a surjection
        \begin{equation}
            \widetilde{\psi}:k[V]_{\overline{g}}\rightarrow \mathcal{O}_V(D_V(g))
        \end{equation}
        where $\psi: k[V]\rightarrow \mathcal{O}_V(D_V(g))$ is the map that takes a polynomial 
        and restricts it to $D_V(g)$ and $\widetilde{\psi}$ is the corresponding map induced by the universal
        property of localization. Furthermore, if $V$ is irreducible, this map is an isomorphism. 
    \end{lemma}
\end{dang}
Proof: First we have $D_V(g)=\bigcup_{i=1}^\ell D_V(g_i)$ and since $\varphi$ is regular, there exists 
$f_i, h_i\in k[\underline{X}],\;i=1,\dots,\ell$ such that 
\begin{equation}
    q\in D_V(g_i)\Rightarrow h_i(q)\neq 0\;(\text{so }q\in D_V(h_i)),\;\varphi\big|_{D_V(g_i)}=\frac{f_i}{h_i}\bigg|_{D_V(g_i)}.
\end{equation}
Now, we have two fundamental relations:
\begin{equation}
    \left\{\begin{aligned}
        &D_V(g_i)\subset D_V(h_i),\;\forall i=1,\dots,\ell\\
        &D_V(g)=\bigcup_{i=1}^\ell D_V(g_i)
    \end{aligned}\right.
\end{equation}
Using $V=V(F_1,\dots, F_m)$ and Hilbert Nullstelensatz~(\ref{theorem : Hilbert Nullstelensatz}), we get 
\begin{equation}
    \left\{\begin{aligned}
        &\sqrt{(F_1,\dots,F_m,g_i)}\subset \sqrt{(F_1,\dots,F_m,h_i)},\;\forall i=1,\dots,\ell\quad (*)\\
        &\sqrt{(F_1,\dots,F_m,g)}=\sqrt{(F_1,\dots,F_m,g_1,\dots,g_\ell)}\quad(**)
    \end{aligned}\right.
\end{equation}
From $(*)$, we deduce that $g_i^{n_i}=\sum_j A_jF_j+B_ih_i$ and restricting to $V$ we get 
$g_i^{n_i}=B_ih_i$. Now, since $\varphi\big|_{D_V(g_i)}=\frac{B_i f_i}{g_i^{n_i}}$ and 
$D_V(g_i)=D_V(g_i^{n_i})$, we can \textcolor{red}{change notation} $B_i f_i\to f_i$ and $g_i^{n_i}\to g_i$.\par 

Now we do a trick: since $g_jf_i - f_jg_i = 0$ on $D_V(g_i)\cap D_V(g_j)=D_V(g_ig_j)$, we get 
\begin{equation}
    g_ig_j(g_jf_i-f_jg_i) = 0\;\text{ on }V.
\end{equation}
Furthermore, we still have $\varphi\big|_{D_V(g_i)}=g_if_i / g_i^2$.
So we \textcolor{red}{change notation} again: $g_i^2\to g_i$ and $g_if_i\to f_i$, so we have (on $V$)
\begin{equation}
    g_if_j - g_jf_i = 0,\;\forall i,j=1,\dots,\ell\quad (***).
\end{equation}
Now, by equation $(**)$, there exist $m>0$ such that 
\begin{equation}
    g^m = \sum_jA_jF_j + \sum_kB_kg_k 
\end{equation}
and we get on $V$:
\begin{equation}
    g^m = \sum_k B_kg_k.
\end{equation}
Now, by using equation $(***)$, we have (on $V$)
\begin{equation}
    f_ig^m = \sum_kB_kf_ig_k = \sum_kB_kf_kg_i = (\sum_kB_kf_k)g_i.
\end{equation}
So if $q\in D_V(g_i)$, we have $\varphi(q) = \tfrac{f_i}{g_i}(q)=\tfrac{f}{g^m}(q)$, where we set $f=\sum_kB_kf_k$.\par 
This proves that $\widehat{\psi}$ is surjective. To prove injectivity, we observe that 
$f/g^m|_{D_V(g)}=0\iff f|_{D_V(g)}=0\iff fg |_V=0$. So if $V$ is irreducible, then $fg\in I(Y)\Rightarrow f\in I(Y)$
because $I(Y)$ is prime, so $\overline{f}=0$. \qed

\begin{dang}
    \begin{corollary}
        Since $V=D_V(1)$, we have 
        \begin{equation}
            \mathcal{O}_V(V)\cong k[V]_{(1)}\cong k[V].
        \end{equation}
    \end{corollary}
\end{dang}

\begin{dang}
    \begin{example}
    Normally, it's not obvious how to describe the set of units $\callO_V(D_V(G))^*$ even if $V$ is a 
    irreducible affine algebraic set. For example, if $V=k^1$, then $k[V]=k[X]$. So, for $G=X(X-1)$ we have 
    $\callO_V(D_V(G))^*\cong k^*\langle X^m, (X-1)^m: m\in\bbZ\rangle$, by the identification
    \begin{equation}
        \callO_V(D_V(G))\cong k[X]_{X(X-1)}\cong k[X,X^{-1},(X-1)^{-1}].
    \end{equation}
\end{example}
\end{dang}



\begin{dang}
    \begin{definition}
        A morphism between two affine algebraic sets $V_1\subset k^n$ and $V_2\subset k^m$ is a function 
        $\varphi:V_1\to V_2$ such that for any open set $U\subset V_2$, we have 
        \begin{equation}
            \varphi_U^{\#}:\mathcal{F}(U,k)\overset{-\circ\varphi}{\rightarrow}\mathcal{F}(\varphi^{-1}(U),k),\;
            \varphi_U^{\#}(\mathcal{O}_{V_1}(U))\subset \mathcal{O}_{V_2}(\varphi^{-1}(U)).
        \end{equation}
    \end{definition}
\end{dang}

Using more fancing words, the previous definition says that a continuous map $\varphi:V_1\to V_2$ is a morphism 
if the composition map $\varphi_U^{\#}:\mathcal{F}(U,k)\overset{-\circ\varphi}{\rightarrow}\mathcal{F}(\varphi^{-1}(U),k)$
induces a sheaf morphism 
\begin{equation}
    \varphi^\#(U):\mathcal{O}_{V_2}\to f_*\mathcal{O}_{V_1}.
\end{equation}
We use the notation $f_*\mathcal{O}_{V_1}$ for direct image sheaf, it's defined by 
$f_*\mathcal{O}_{V_1}(U)=\mathcal{O}_{V_1}(f^{-1}(U))$.


\section{Lecture 4}


\begin{dang}
    \begin{definition}
        A \underline{function space} is a topological space $X$ with a sheaf $\mathcal{O}_X$ such that 
        for any open set $U\subset X$, we have 
        \begin{itemize}
            \item[(1)] $\mathcal{O}_X(U)\overset{\text{subring}}{\subset} \{U\overset{\text{continuous}}{\to}k\}$
            \item[(2)] Let $f\in\mathcal{O}_X(U)$, then $f$ has an inverse in $\mathcal{O}_X(U)$ if and only if 
            $f(q)\neq0,\;\forall q\in U$.
        \end{itemize}
    \end{definition}
\end{dang}


\begin{dang}
\begin{lemma}
    An affine algebraic set $V\subset k^n$ equiped if the sheaf of regular functions is a function space.
\end{lemma}
\end{dang}
Proof: Let $U\subset X$ be an open set and $f\in\mathcal{O}_X(U)$. Then, by definition, there exist an open cover 
$U=\bigcup_\lambda U_\lambda$ such that $f|_{U_\lambda}=f_\lambda/g_\lambda$, where $g_\lambda(q)\neq0\;\forall q\in U_\lambda$. 
To prove that $f$ is continuous, it's sufficient to prove that $f^{-1}(a)\subset V$ is closed for all $a\in k$. 
So we prove that $(f|_{U_\lambda})^{-1}(a)$ is closed in $V$ for all $\lambda$, which can be verified by the 
equality
\begin{equation}
    V\cap V(f_\lambda - ag_\lambda) = (f|_{U_\lambda})^{-1}(a).
\end{equation}
The second item is clear.\qed 

\begin{dang}
    \begin{definition}
        Let $(X,\mathcal{O}_X)$ and $(Y,\mathcal{O}_Y)$ be two function spaces. We define a morphism $X\to Y$ to be a 
        a continuous map $f:X\to Y$ which induces a sheaf map $f^\#:\mathcal{O}_Y\to f_*\mathcal{O}_Y$ in the sence that 
        if $f^\#(U)(g)=g\circ f$ for $g\in\mathcal{O}_Y(U)$, then $f^\#(U)(\mathcal{O}_Y(U))\subset \mathcal{O}_X(f^{-1}(U))$.
    \end{definition}
\end{dang}

\begin{remark}
    Being a morphism of function spaces is a local property. In fact, let $f:X\to Y$ be a morphism of function spaces and 
    suppose that $Y=\bigcup_\lambda V_\lambda$ and for each $\lambda$:
    \begin{equation}
        f|_{f^{-1}(V_\lambda)}:f^{-1}(V_\lambda)\to V_\lambda
    \end{equation}
    is a morphism. Then $f$ is a morphism.
\end{remark}

\begin{dang}
    \begin{definition}
        An \underline{affine algebraic variety} is a function space $(X,\mathcal{O}_X)$ isomorphic to 
        an affine algebraic set with Zariski topology and sheaf of regular functions.
    \end{definition}
\end{dang}


\begin{dang}
    \begin{definition}
        Let $(X,\mathcal{O}_X)$ be a function space, and $U\subset X$ be an open set with induced topology and 
        $\mathcal{O}_U=\mathcal{O}_X|_U$. Then $(U,\mathcal{O}_U)$ is a function space called 
        \underline{induced open subspace}.
    \end{definition}
\end{dang}

\begin{dang}
    \begin{definition}
        An \underline{abstract algebraic variety} is a function space $(X,\mathcal{O}_X)$ such that there exist an open covering  
        $X=\bigcup_\lambda U_\lambda$ with the property that for each $\lambda$, the induced 
        open subspace $(U_\lambda,\mathcal{O}_{U_\lambda})$ is an affine algebraic variety.
    \end{definition}
\end{dang}

\begin{dang}
    \begin{remark}
    In many texts we have the additional hypothesis of irredutibility in the definition of algebraic varieties.
\end{remark}
\end{dang}

\begin{dang}
    \begin{definition}
        An affine algebraic variety is said to be \underline{quasi-affine} if it's isomorphic, as a function space, 
        to an induced open subspace of an affine algebraic variety.
    \end{definition}
\end{dang}

\begin{dang}
    \begin{lemma}
        Let $(X,\mathcal{O}_X)$ be a function space and $Y\subset k^n$ be a quasi affine algebraic set
    \end{lemma}
\end{dang}


\begin{dang}
    \begin{lemma}\label{lemma: coordinate criteria to a map be a morphism of function spaces}
        Let $(X,\mathcal{O}_X)$ be a function space and $Y\subset k^n$ be a quasi-affine algebraic set. 
        Then $f:X\to Y$ is a morphism if and only $f=(g_1,\dots,g_n)$ with $g_i\in\mathcal{O}_X(X)$. Here, 
        we think $Y$ as an induced open subspace of an affine algebraic set.
    \end{lemma}
\end{dang}
Proof: $(\Rightarrow)$ The coordinate functions $x_i|_Y:Y\to k$ are regular functions, by definition. Now 
$g_i = x_i\circ f = f^\#(Y)(x_i)\in\mathcal{O}_X(f^{-1}(Y))=\mathcal{O}_X(X)$, because 
$x_i\in\mathcal{O}_Y(Y)=\{Y\overset{\text{regular}}{\to}k\}$.\par 
$(\Leftarrow)$ A closed set of $Y$ is of the form $V(F)\cap Y$ 
\textcolor{red}{(would not be of the form $V(F_1,\dots, F_m)$ in general ?)}. However
$f^{-1}(V(F)\cap Y)=\{p\in X: F(g_1,\dots, g_n)(p)=0\}$ which is closed in $X$ provided we show that 
$F(g_1,\dots, g_n)$ is continuous. However, since $\mathcal{O}_X(X)$ is a ring, we have $F(g_1,\dots, g_n)\in\mathcal{O}_X(X)$. 
Furthermore, if $U\subset Y$ is an open set and $\varphi:U\to k$ is regular, we need to show that 
$f^{\#}(U)(\varphi)\in\mathcal{O}_X(f^{-1}(U))$. Let $U^\prime\subset U$ be such that 
$\varphi|_{U^\prime}=F/G,\;G(q)\neq0\;\forall q\in U^\prime$. Then we have 
\begin{equation}
    f^\#(U^\prime)(\varphi)=\varphi\circ f|_{f^{-1}(U^\prime)} = 
    \frac{F(g_1,\dots, g_n)}{G(g_1,\dots, g_n)}\bigg|_{f^{-1}(U^\prime)}.
\end{equation}
Furthermore 
\begin{equation}
    G\circ f(q) = G(g_1,\dots, g_n)(q) \neq 0,\;\forall q\in f^{-1}(U^\prime).
\end{equation}
Since $G(g_1,\dots,g_n)|_{f^{-1}(U^\prime)}\in\mathcal{O}_X(f^{-1}(U^\prime))$, the previous equation 
shows that the inverse $1/G(g_1,\dots,g_n)$ also is an element of $\mathcal{O}_X(f^{-1}(U^\prime))$. Now, let 
$U=\bigcup_\lambda U_\lambda,\;\varphi|_{U_\lambda}=F_\lambda/G_\lambda$, then we see that 
\begin{equation}
    f^\#(U_\lambda)(\varphi|_{U_\lambda})\subset \mathcal{O}_X(f^{-1}(U_\lambda)).
\end{equation}
The glueing axiom shows that $f^\#(U)(\varphi)\in\mathcal{O}_X(f^{-1}(U))$.\qed


\begin{dang}
    \begin{lemma}
        Let $V\subset k^n$ be an affine algebraic set and $D_V(F)\subset V$ an affine algebraic set. Then 
        the induced open subspace 
        $(D_V(F), \mathcal{O}_{D_V(F)})$ is a quasi-affine algebraic variety, by definition. However, we have more, 
        it's also an affine algebraic variety.
    \end{lemma}
\end{dang}
Proof: Let $V=V(G_1,\dots, G_m)$ and define the map 
\begin{equation}
    \varphi: D_V(F)\rightarrow V(x_0F-1, G_1,\dots, G_m)\subset k^{n+1}:(a_1,\dots,a_n)\mapsto 
    (\frac{1}{F(a_1,\dots, a_n)},a_1,\dots, a_n)
\end{equation}
By the lemma~(\ref{lemma: coordinate criteria to a map be a morphism of function spaces}), we see that 
$\varphi$ is a morphism. Furthermore, it has a set theoretical inverse given by the projection 
\begin{equation}
    \psi:V(x_0F-1, G_1,\dots, G_m)\to D_V(F):(a_0,\dots,a_n)\mapsto (a_1,\dots, a_n).
\end{equation}
By the same lemma, this map is also a morphism. Since $f\psi = \mathds{1}_{D_V(x0F-1,\dots)}$ and 
$\psi f = \mathds{1}_{D_V(F)}$, we see that $\varphi$ is an isomorphism (the induced 
sheaf morphism would be an isomorphism by construction).\qed

\begin{dang}
    \begin{remark}
        Let $(X,\mathcal{O}_X)$ be an affine algebraic variety and 
        $\varphi:(V,\mathcal{O}_V)\to (X,\mathcal{O}_X)$ an isomorphism, where $V\subset k^n$ is an algebraic set. 
        Then this isomorphism, induces another isomorphism
        \begin{equation}
            (D_V(F),\mathcal{O}_{D_V(F)}) \to (X_f,\mathcal{O}_{X_f})
        \end{equation}
        where $F=f\circ\varphi$ and $X_f = \varphi(D_V(F))=\{p\in X: f(p)\neq0\}$.
    \end{remark}
\end{dang}

\begin{dang}
    \begin{corollary}
        A quasi-affine algebraic variety is an abstract algebraic variety.
    \end{corollary}
\end{dang}
Proof: A quasi-affine algebraic variety has a topological basis consisting of principal open subsets, which 
are affine algebraic varieties.



\section{Lecture 5 and 6}

\begin{dang}
    \begin{theorem}\label{theorem: glueing theorem}
        Suppose that we have the data:
\begin{itemize}
    \item $(X_\lambda)_{\lambda\in\Lambda}$ function spaces over $k$
    \item $\forall\lambda,\mu\in\Lambda$ let $U_{\lambda\mu}\subset X_\lambda$
    \item $\forall\lambda,\mu\in\Lambda$ an isomorphism
    \begin{equation}
        \varphi_{\lambda\mu}:U_{\lambda\mu}\rightarrow U_{\mu\lambda}
    \end{equation} 
\end{itemize}
Furthermore, assume that 
\begin{itemize}
    \item $U_{\lambda\lambda} = X_\lambda,\;\varphi_{\lambda\lambda} = \mathds{1}_{X_\lambda}$
    \item $\forall\lambda,\mu,\nu\in\Lambda$, we have 
    \begin{equation}
        \begin{split}
            &\varphi_{\lambda\mu}(U_{\lambda\mu}\cap U_{\lambda\nu})\subset U_{\mu\lambda}\cap U_{\mu\nu}\\
            &\varphi_{\mu\nu}\varphi_{\lambda\mu}\big|_{U_{\lambda\mu}\cap U_{\lambda\nu}}=\varphi_{\lambda\nu}\big|_{U_{\lambda\mu}\cap U_{\lambda\nu}}
        \end{split}
    \end{equation}
\end{itemize}
Then, $X=\coprod_{\lambda\in\Lambda}X_{\lambda}/\sim$, where we put $(\lambda,x)\sim (\mu,x^\prime)$ if and only if 
$\varphi_{\lambda\mu}(x)=x^\prime$. One can verify that this defines an equivalence relation. Furthermore, 
$X$ has a topology defined as the small topology such that all the maps $i_\lambda:X_\lambda\to X:x\mapsto [(\lambda,x)]$ 
are continuous. We define a sheaf 
\begin{equation}
    \mathcal{O}_X(U)=\{f:U\to k: fi_\lambda\in \mathcal{O}_{X_\lambda}(i_\lambda^{-1}(U)),\;\forall\lambda\in\Lambda\}.
\end{equation}
With this sheaf, $(X,\mathcal{O}_X)$ is a function space satisfying the following universal property: 
for any function space $Y$ and morphisms $j_\lambda: X_\lambda\to Y$such that:
\begin{equation}
    j_\mu \varphi_{\lambda\mu} = j_\lambda|_{U_{\lambda\mu}},
\end{equation}
there exist an unique morphism $f: X\to Y$ such that 
$fi_\lambda = j_\lambda$ (morphisms here, means morphisms of function spaces !). Furtherore, since 
$i_\mu^{-1}i_\lambda(X_\lambda)=U_{\mu\lambda}\subset X_\mu,\;\forall\mu$, then $i_\lambda(U_\lambda)$ is open and 
$i_\lambda$ are isomorphisms onto its images. So, if each $X_\lambda$ is an affine algebrai variety, then 
$X$ is an abstract variety.
    \end{theorem}
\end{dang}
Proof: Tedious.\qed

\begin{dang}
    \begin{definition}
        An abstract variety $X$ is separable if for any pair of affine open sets $U,V\subset X$, we have 
        \begin{itemize}
            \item[(1)] $U\cap V$ is affine
            \item[(2)] The induced map 
            \begin{equation}
                \mathcal{O}(U)\otimes\mathcal{O}(V)\rightarrow \mathcal{O}(U\cap V): f\otimes g\mapsto f|_{U\cap V}\cdot 
                g|_{U\cap V}
            \end{equation}
            is surjective.
        \end{itemize}
    \end{definition}
\end{dang}


\begin{example}
    Let $X_0= X_1 = k^1$ and $U_{01}=U_{10}=k^1-\{0\}$. Define a map $\varphi_{01}:U_{01}\to U_{10}:x\mapsto 1/x$. 
    We claim that $X = \mathbb{P}^1$ in the variety obtained by glueing $X_0$ and $X_1$ through $\varphi_{01}$. 
    In fact, for $i=0,1$, let $g_i: X_i\to Y$ be two morphisms such that 
    \begin{equation}
        g_1\varphi_{01}=g_0\big|_{U_{01}}.
    \end{equation}
    That is, $g_1(1/x)=g_0(x)$. Then the following map is well defined:
    \begin{equation}
        \begin{split}
             g:&X\to Y\\
             &[1:y]\mapsto g_0(y)\\
             &[x:1]\mapsto g_1(x)
        \end{split}
    \end{equation}
    Moreover, by construction we have $g|_{U_i}=g_i\varphi_i$, where $\varphi_i: U_i\to k^1$ 
    (here, $U_i=\{[x_0,x_1]:x_i\neq0\}$) are the standard charts of 
    the projective plane $X$. So we se that $g$ is locally a morphism, i.e., it's a morphism.
\end{example}

\section{Lecture 6,7,8}
\begin{remark}
    Take the product topology on $k^1\times k^1$, where $k^1$ is equiped with Zariski topology. Let 
    $U=k^1-S$ and $V=k^1-T$ ($S,\;T$ finite) be two open sets. Then 
    $U\times V = k^1\times k^1 - (S\times k^1\cup k^1\times T)$ and in general, for a familiy of open sets 
    $U_\lambda, V_\lambda$, we have the equality.
    \begin{equation}
    \begin{split}
    &(\bigcup_\lambda U_\lambda\times V_\lambda)^c = \bigcap_\lambda S_\lambda\times k^1\cup k^1\times T_\lambda\\
    &\quad=(\bigcap_\lambda S_\lambda\times k^1)\cup (k^1\times \bigcap_\lambda T_\lambda)    
    \end{split}
    \end{equation}
    where $S_\lambda, T_\lambda$ are all finite. So any closed set $F\subset k^1\times k^1$ is of the form 
    \begin{equation}
        F = S\times k^1 \cup k^1\times T,\; \#S,\#T<\infty
    \end{equation}
    in the product topology. In particular, if $\# k =\infty$, then $V(X,Y)=\{0\}$ is not closed in the product topology.
\end{remark}

\TODO remains to study lect 8, Nullstelensatz for projective spaces and ``universal compactness'' theorem.
\subsection{Projective varieties}
We define $\bbP^m = (k^{m+1}\setminus\{0\}) \big/\sim$. As usual, define charts 
$\varphi_i: D_+(X_i)\to k^{m}: [x_0,\dots,x_m]\mapsto 
(\frac{x_0}{x_i},\dots,\frac{x_{i-1}}{x_i},\frac{x_{i+1}}{x_i},\dots,\frac{x_m}{x_i})$. Then we see that the transition maps 
$\varphi_i\varphi_j^{-1}: D(X_i)\subset k^m\to D(X_j)\subset k^m$ are isomorphisms of affine varieties. 
Furthermore, $\bbP^m$ has a topology induced by the open affine covering $\bbP^m =\bigcup_{i=0}^mD_+(X_i)$ where we set 
$U\subset\bbP^m$ to be open if $U\cap D_+(X_i)\subset D_+(X_i)$ is open $\forall i$. With this definition it's clear 
that $D_+(X_i)$ is open, because $\varphi_i(D_+(X_iX_j))= D(X_j)\subset k^m$ is open $\forall j$. The sheaf of regular 
function is defined by using the charts:
\begin{equation}
    \mathcal{O}_{\bbP^m}(U)=\{f:U\to k: f\varphi_i^{-1}\in\callO_{k^m}(\varphi_i(U)),\;\forall i\}
\end{equation}
In this condition, we see that $\bbP^m$ is an abstract variety.

\begin{dang}
    \begin{proposition}
        Let $\varphi_{ij}=\varphi_j\varphi_i^{-1}:\varphi_i(D(X_iX_j))\to\varphi_j(D(X_iX_j))$. Write 
        $U_{ij} = \varphi_i(D(X_iX_j))\subset k^m$. Then $\varphi_{ij}:U_{ij}\to U_{ji}$ satisfies all the conditions 
        of theorem~(\ref{theorem: glueing theorem}). We claim that $\bbP^m$ is isomorphic to the variety obtatined 
        by glueing all $k^m$ through the isomorphisms $\varphi_{ij}$. 
    \end{proposition}
\end{dang}
Proof: We have a well defined bijection 
\begin{equation}
    \bbP^m\rightarrow X=\coprod_{i=0}^m\{i\}\times k^{m}\bigg/\sim\;:D_+(X_i)\ni[x_0,\dots,x_m]\mapsto [\varphi_i([x_0,\dots,x_m])_i]
\end{equation}
where we write $(\underline{x})_i$ for an element of $\{i\}\times k^m$. It remains to show that this map is an isomorphism. 
However, at the level of charts we have the identity map. \qed

\subsection{Projective Nullstelensatz}
\begin{dang}
    \begin{proposition}
        Let $I\subset k[X_0,\dots, X_m]$ be a ideal generated by homogeneous polynomial. Then 
        $F = \sum_i F_i\in I\iff F_i\in I\;\forall i$ where $F_i$ are the homogeneous components of $F$. In each 
        case we say that the ideal $I$ is \textit{homogeneous}.
    \end{proposition}
\end{dang}

\begin{dang}
    \begin{proposition}
        An ideal $I$ is homogeneous if and only if $\sqrt{I}$ is homogeneous.
    \end{proposition}
\end{dang}

\begin{notation}
    Let $I\subset k[X_0,\dots,X_m]$ be a homogeneous ideal, then we denote 
    \begin{equation}
        V_p(I) = \{[x_0,\dots,x_m]\in\bbP^m: F(x_0,\dots,x_m)=0,\;\forall F\in I\text{ homogenous}\}.
    \end{equation}
    Furthermore, if $V\subset \bbP^m$ is a subset, we write 
    \begin{equation}
        I_p(V)=\langle S\rangle,\;S=\{F\in k[X_0,\dots,X_m]: F\text{ is homogeneous and }F(p)=0,\;\forall p\in V\}.
    \end{equation}
\end{notation}


\begin{dang}
    \begin{theorem}\label{theorem: projective Nullstelensatz}
        Let $I\subset k[X_0,\dots,X_m]$ be a homogeneous ideal. Then 
        \begin{itemize}
            \item[(1)] $V_p(I)=\varnothing\iff \sqrt{I}\supset (X_0,\dots,X_m)$
            \item[(2)] If $V_p(I)\neq\varnothing$, then $I_pV_p(I)=\sqrt{I}$ 
        \end{itemize}
    \end{theorem}
\end{dang}
Proof: For the first statement, we observe that 
$V_p(I)=\varnothing\iff V_a(I)=\{0\}$. However, by the Nullstelensatz we have 
\begin{equation}
    \sqrt{I} = I_aV_a(I)=I_a(\{0\})=(X_0,\dots,X_m).
\end{equation}
For the second statement, we observe that if $V_p(I)\neq\varnothing$, then 
\begin{equation}
    \begin{split}
        I_pV_p(I)&=\langle F\text{ homogeneous }: F[x_0,\dots,x_m]=0,\;\forall [x_0,\dots,x_m]\in V_p(I)\rangle\\
         &=\langle F\text{ homogeneous }: F[x_0,\dots,x_m]=0,\;\forall (x_0,\dots,x_m)\in V_a(I)\setminus\{0\}\rangle\\
         &=\langle F\text{ homogeneous }: F(x_0,\dots,x_m)=0,\;\forall (x_0,\dots,x_m)\in V_a(I) \rangle\\
         &=\langle F\text{ homogeneous }: F\in I_aV_a(I)\rangle\\
         &=\langle F\text{ homogeneous }: F\in \sqrt{I}\rangle = \sqrt{I}.
    \end{split}
\end{equation}
The last line follows because $\sqrt{I}$ is a homogenous ideal.\qed

\subsection{Segre embedding}
\begin{remark}
    Here, to ease notation, we will write $x=[x_0,\dots, x_m]$ sometimes.
\end{remark}
The Segre embedding is the map 
\begin{equation}
    \begin{split}
          &s:\bbP^{m}\times\bbP^n\rightarrow \bbP^{(m+1)(n+1)-1}\\
          &\quad([x_0,\dots,x_m],[y_0,\dots, y_n])\mapsto 
          \left[\begin{matrix}
            x_0y_0&x_0y_1&\dots&x_0y_n\\
            x_1y_0&x_1y_1&\dots&x_1y_n\\
            \vdots&\vdots&\vdots&\vdots\\
            x_my_0&x_my_1&\dots&x_my_n
          \end{matrix}\right]
    \end{split}
\end{equation}
where we parametrize $\bbP^{(m+1)(n+1)-1}$ with projective coordinates $z_{ij},\;0\leq i\leq m,0\leq j\leq n$.
We claim that $s(\bbP^m\times\bbP^n)$ is a closed subset. In fact, if we let 
$z_{ij}=x_iy_j$, then we have $z_{ij}z_{pq}=z_{iq}z_{pj}$ and 
\begin{equation}
    s(\bbP^m\times\bbP^n)\subset V=V_p(z_{ij}z_{pq}-z_{pj}z_{iq},\;\forall p,q,i,j)\quad(*).
\end{equation}
Conversely, if we take $z\in V$, then $z_{ij}\neq0$ for some $i,j$. So we assume $z_{ij}=1$. Then we see that 
$z_{pq} = z_{pj}z_{iq}$, so we can take $x_p  = z_{pj}$ and $y_q=z_{iq}$ to conclude that 
$s(x,y)=z$. Then we conclude the equality of $(*)$. Furthermore, we claim that this map is injective. Let 
$s(x,y)=s(x^\prime,y^\prime)$, then $x_py_q=\lambda x_p^\prime y_q^\prime$. Fix $q$ such that $y_q\neq0$, then 
$x_p = (\frac{\lambda y_q^\prime}{y_q})x_p^\prime,\;\forall p$ and $x=x^\prime$. In a similar way we see that $y=y^\prime$.\par 

As a corollary, we see that $s$ mapst $\bbP^m\times\bbP^n$ bijectively to a closed subset $V\subset\bbP^{(m+1)(n+1)-1}$. 
We put a structure of abstract variety on $\bbP^m\times \bbP^n$ in such a way that $s$ is an isomorphism onto $V$. So we 
see for example, that 
\begin{equation}
    \callO_{\bbP^m\times\bbP^n}(U)=(-\circ s)(\callO_V(s(U))).
\end{equation}
In a concrete way, a section $s\in \callO_{\bbP^m\times\bbP^n}$ is a quotient of bihomogeneous polynomials of bi degree $(d,d)$
each one, for some $d\geq0$. In fact, if $F(z)\in k[\underline{z}]$ is homogeneous, then 
\begin{equation}
    F(z)=\sum a_\alpha\prod_{i,j}z_{ij}^{\alpha_{ij}},\; (i,j)\in\{0,\dots,m\}\times\{0,\dots,n\}.
\end{equation}
Therefore 
\begin{equation}
    \begin{split}
        F\circ s (x,y) &= \sum a_\alpha\prod_{i,j}x_i^{\alpha_{ij}}y_j^{\alpha_{ij}}\\
             &=\sum a_\alpha(\prod_i x_i^{\sum_j\alpha_{ij}})(\prod_j y_j^{\sum_i\alpha_{ij}})
    \end{split}
\end{equation}
where $\sum_i \sum_j \alpha_{ij}=\sum_j\sum_i\alpha_{ij}=|\alpha|=d$, the degree of $F$. It's easy to see that a quotient of
bihomogeneous polynomials of bidegree $(d,e)$ each one, 
is also a section of $\callO_{\bbP^m\times\bbP^n}$. In fact, let $G/H$ be a quotient of two bihomogeneous polynomials
of bidegree $(d,e)$ each one. Then, in $D(x_i)\times D(y_j)$ we have 
\begin{equation}
    \frac{G}{H} = \frac{G x_i^ey_j^d}{Hx_i^ey_j^d}.
\end{equation}
So we can assume both $G,H$ to be of the same bidegree $(d+e,d+e)$. Moreover, let $F$ be a bihomogeneous polynomial 
of bidegree $(d,d)$, then 
\begin{equation}
    \begin{split}
        F(x,y) &= \sum a_{\alpha\beta}x^\alpha y^\beta\\
        &=\sum a_{\alpha\beta} x_{i_1} y_{j_1}\cdots x_{i_d}y_{j_d}\; (1\leq i_k\leq m, 1\leq j_k\leq n)\\
        &=\sum a_{\alpha\beta} z_{i_1j_1}\cdots z_{i_dj_d}
    \end{split}
\end{equation}
is homogeneous of degree $d$.

\begin{dang}
    \begin{lemma}
        The variety $\bbP^m\times \bbP^n$ constructed via the Segre embedding is a product.
    \end{lemma}
\end{dang}
Proof: Let $f:Z\to \bbP^m$ and $g:Z\to \bbP^n$ be morphisms. 
We draw the following diagram
\begin{center}
    \begin{tikzcd}
       &D(x_i)\times D(y_j)\arrow[r,"s"]\arrow[dd,"{(\varphi_i,\psi_j)}"]& D(z_{ij})\cap V&\\
     U= f^{-1}(D(x_i))\cap g^{-1}(D(y_j))\arrow[ru,"{(f,g)|_U}"]\arrow[rd,swap,"{(\varphi_i f|_U,\psi_j g|_U)}"] &&&\\
        &k^m \times k^n\arrow[ruu,swap,dashed,"{s\circ (\varphi_i^{-1},\psi_j^{-1})}"]&&
    \end{tikzcd}
\end{center} 
We see that $\varphi_if\big|_U, \psi_jg\big|_U$ are morphisms, so $(\varphi_i f|_U,\psi_j g|_U)$ is a 
morphism because it's a mapping into a product of affine varieties and we can apply the affine criteria. So, if we show that 
$s\circ (\varphi_i^{-1},\psi_j^{-1})$ is a morphism, we are done. However, we see that 
\begin{equation}
    \xi_{ij}\circ s \circ(\varphi_i^{-1},\psi_j^{-1}):k^m\times k^n\rightarrow k^{(m+1)(n+1)}:
    (x,y)\mapsto (\dots, x_i,\dots, y_j,\dots, x_iy_j,\dots)
\end{equation}
which is clearly a morphism.\qed 

\begin{remark}
    The above proof also shows that $(\varphi_i,\psi_j):D(x_i)\times D(y_j)\to k^m\times k^n$ is a morphism. Furthermore, 
    let $X\subset\bbP^m, Y\subset\bbP^n$ be two closed subsets. Then, is 
    $X\cap D(x_i) \times Y\cap D(y_j)\subset D(x_i)\times D(y_j)$ closed ?. We have that 
    $\varphi_i(X\cap D(x_i))\subset k^m, \psi_j(Y\cap D(y_j))\subset k^n$ are closed. Therefore, the product
    \begin{equation}
        \varphi_i(X\cap D(x_i))\times \psi_j(Y\cap D(y_j)) = (\varphi_i,\psi_j)(X\cap D(x_i) \times Y\cap D(y_j))
    \end{equation}
    is closed in $k^m\times k^n$. Here we use the standard observation that $V(S)\times V(T) = V(S\cup T)$, where 
    $S\subset k[\underline{X}], T\subset k[\underline{y}]$ with both rings sitting inside $k[\underline{x},\underline{y}]$. 
    Since $(\varphi_i,\psi_j)$ is a homeomorphism, this proves that the answer is positive.
\end{remark}

\begin{dang}
    \begin{corollary}\label{corollary: product of closed subsets in the product of projective varieties}
        Let $X\subset \bbP^m, Y\subset \bbP^n$ be two closed subsets. Then 
        $X\times Y\subset \bbP^m\times \bbP^n$ is closed.
    \end{corollary}
\end{dang}

\begin{remark}
    The product of the affine spaces $k^m$ and $k^n$ if $k^{m}\times k^n\cong k^{m+n}$. So, if 
    $f:Z\to k^m,g:Z\to k^n$ are morphisms, then we draw the diagram.
    \begin{center}
        \begin{tikzcd}
            &k^m\times k^n\arrow[dd,"p_1"]\arrow[ddr,"p_2"]&\\
          Z\arrow[ru,dashed,"{(f,g)}"]\arrow[rd,swap,"f"]\arrow[rrd,"g"] &&\\
            &k^m&k^n
        \end{tikzcd}
    \end{center}
    However, $f$ being a morphism means that $x_i\circ f$ and $y_j\circ g$ are morphims for all 
    $i,j\;(1\leq i\leq m,1\leq j\leq n)$, where $x_i$ and $y_j$ are the standard coordinate functions. However, the 
    set of coordinate functions of the product $k^m\times k^n$ is $(x_1,\dots, x_m,y_1,\dots, y_n)$ because of the 
    identification $k^m\times k^n\cong k^{m+n}$. So, again we see that 
    \begin{equation}
        (x_1\circ f,\dots, x_m\circ f,y_1\circ g,\dots, y_n\circ g)=(f,g)
    \end{equation}
    is a morphism. Here we used lemma~(\ref{lemma: coordinate criteria to a map be a morphism of function spaces})
    several times.
\end{remark}
\subsection{Grassmanians and Flag varieties}
Let $\Gr_d(V)$ be the set of $d-$dimensional subspaces of $V$. We have the following embedding
\begin{equation}
    \psi: \Gr_d(V)\hookrightarrow \bbP(\bigwedge^dV): W=\text{Span}(w_1,\dots, w_d)\mapsto [w_1\wedge\cdots\wedge w_d].
\end{equation}
The map is well defined, in fact, if $\{w_i\}_{i=1}^d$ and $\{w^\prime_i\}_{i=1}^d$ are two basis of $W$ and 
$w_i^\prime = \sum_j b_{ij}w_i$, then 
$w_1^\prime\wedge\cdots\wedge w_d^\prime = \det(b_{ij})_{i,j=1}^d w_1\wedge\cdots\wedge w_d$. Now, let 
$v_1,\dots, v_n$ be a basis of $V$ and $u_1,\dots,u_d\in V$ be arbitrary vectors. Write 
$u_i = \sum_ja_{ij}v_j$, then 
\begin{equation}
    \begin{split}
        u_1\wedge\cdots\wedge u_d &= (\sum_{j_1=1}^na_{1j_1}v_{j_1})\wedge\cdots\wedge(\sum_{j_d=1}^na_{dj_d}v_{j_d})\\
        &=\sum_{j_1=1}^n\cdots\sum_{j_d=1}^n a_{1j_1}\cdots a_{dj_d}v_{j_1}\wedge\cdots\wedge v_{j_d}\\
        &=\sum_{j_1<\cdots<j_d}\det(a_{ij_k})_{i,k=1}^d v_{j_1}\wedge\cdots\wedge v_{j_d}
    \end{split}
\end{equation}
So we see that $u_1\wedge\cdots\wedge u_d=0\iff \text{Rank}(a_{ij})<d\iff \{u_1,\dots, u_d\}$ is linearly dependent. So, if 
\begin{equation}
    w_1\wedge\dots\wedge w_d = \lambda w_1^\prime\wedge\dots\wedge w_d^\prime \neq0,
\end{equation}
then $u\wedge w_1\wedge\dots\wedge w_d=0\iff  u\wedge w_1^\prime\wedge\dots\wedge w_d^\prime=0$, i.e., 
$u\in\text{Span}\{w_i\}_{i=1}^d\iff u\in\text{Span}\{w_i^\prime\}_{i=1}^d$. Now, let 
$\mathcal{J}=\{(j_1,\dots, j_d):1\leq j_1<\cdots<j_d\leq n\}$ and if $J=(j_1,\dots, j_d)\in\mathcal{J}$, write 
$v_J = v_{j_1}\wedge\cdots\wedge v_{j_d}$. Then we have the basis $\{v_J\}_{J\in\mathcal{J}}$ of $\bigwedge^dV$ and we 
write $X_J$ for the corresponding coordinate functions.\par 

Let $W\in\Gr_d(V)$, we see that 
$\psi(W)\in D(X_{(1,\dots,d)})$ if and only if $\det(a_{ij})_{i,j=1}^d\neq0$. In this case, by row reduction, we can assume 
that $a_{ij}=\delta_{ij},\;(1\leq i,j\leq d)$, so $W$ has a basis 
consisting of vectors of the form $w_i = v_i+\sum_{j>d}a_{ij}v_j,\;i=1,\dots,d$. Now define $J_{pq}$ as the tuple 
$(j_1,\dots, j_d)$ such that $j_k=k\;(k < p),\;j_{k}=k+1,\;(p\leq k<d)$ and $j_d=q\in\{d+1,\dots,n\}$. With this notation, we 
see that 
\begin{equation}\label{equation: a_pq coefficients}
    a_{pq}=\pm\det(a_{ij_k})_{j,k=1}^d=\pm X_{J_{pq}}(w_1\wedge\cdots\wedge w_d).
\end{equation}
So, for any 
$J=(j_1,\dots,j_d)\in\mathcal{J}$, the expression $\det(a_{ij_k})_{i,k=1}^d$ is a polynomial 
$f_J(X)\in k[X_{J_{pq}}: 1\leq p\leq d, d<q\leq n]\subset k[X_J:J\in\mathcal{J}\setminus\{(1,\dots,d)\}]$ and we see that 
\begin{equation}\label{equation: local description of a Grassmanian}
    \begin{split}
        &\varphi_{(1,\dots,d)}(\text{Im}(\psi)\cap D(X_{(1,\dots,d)}))\\
        &\quad =V(X_J-f_J(X):\forall J=(j_1,\dots,j_d)\in\mathcal{J}\setminus\{(1,\dots,d), J_{pq}: p,q\})
    \end{split}
\end{equation}
which is closed, so $\text{Im}(\psi)\cap D(X_{(1,\dots,d)})$ is closed. With the same argument applied to the other 
principal open subsets $D(X_J)$, we see that $\text{Im}(\psi)$ is closed. It's important to no note that 
the row echelon form of $A=(a_{ij})$ gives an affine coordinate expression for $\psi(W)\in U_{(1,\dots,d)}$:
\begin{equation}
    \varphi_{(1,\dots,d)}\psi(W) = (f_J(a))_{J\in\mathcal{J}\setminus\{(1,\dots,d)\}}
\end{equation}
where $f_J(a)=\det(a_{ij_k})_{i,k=1}^d,\; J=(j_1,\dots,j_d)$.\par 

Now we define the Flag variety $\mathcal{F}(V)$, it's defined as the $n$-tuple of subspaces  
\begin{equation}
    (F_1,\dots, F_n)\in \Gr_1(V)\times\cdots\times \Gr_n(V)
\end{equation}
such that $F_i\subset F_{i+1},\;1\leq i<n$. It's a closed subset, so it's a projective variety in the above product. 
We prove this in two steps. The first one is to prove that 
\begin{equation}
    \{(F,F^\prime):F\subset F^\prime\}\subset \Gr_d(V)\times\Gr_{d+1}(V)
\end{equation}
is a closed subset. As before, we simplify the problem by showing that the intersection 
\begin{equation}
     \{(F,F^\prime):F\subset F^\prime\}\cap U_{(1,\dots,d)}\times U_{(1,\dots,d,d+1)}\quad(*)
\end{equation}
is closed. So let $F$ and $F^\prime$ be represented by matrices $(a_{ij})\in\text{Mat}_{d,n}(k)$ and 
$(b_{pq})\in\text{Mat}_{d+1,n}(k)$ in row echelon form, i.e., 
$a_{ij}=\delta_{ij},\;j\leq d$ and $b_{pq}=\delta_{pq},\;q\leq d+1$. In this condition 
$F\subset F^\prime$ if and only if the matrix 
\begin{equation}
   M= \begin{pmatrix}
        b_{11}&\cdots& b_{1n}\\
        \vdots&\cdots&\vdots\\
        b_{d+1,1}&\cdots&b_{d+1,n}\\
        a_{11}&\cdots&a_{1n}\\
        \vdots&\cdots&\vdots\\
        a_{d1}&\cdots&a_{dn}
    \end{pmatrix}
\end{equation}
has rank $d+1$. In particular, this means that any $d+2$ minor of $M$ must vanish, which is a polynomial condition 
on the coefficients $a_{ij}, b_{pq}$. By equation~(\ref{equation: a_pq coefficients}), we see that this 
defines a polynomial condition in the affine coordinate system of $\varphi_{(1,\dots,d)}\times\varphi_{(1,\dots,d+1)}$, hence 
intersection $(*)$ is closed.\par 

Now, for the second step, we use corollary~(\ref{corollary: product of closed subsets in the product of projective varieties}) 
to conclude that 
\begin{equation}
    \begin{split}
&S_{d} = \Gr_1(V)\times\cdots\times\Gr_{d-1}(V)\times\{(F,F^\prime)\in\Gr_d(V)\times \Gr_{d+1}(V):F\subset F^\prime\}\times\cdots\\
&\quad\quad\quad\quad\cdots\times\Gr_{d+2}(V)\times\cdots\times \Gr_n(V) 
    \end{split}
\end{equation}
is closed. Therefore 
\begin{equation}
    \mathcal{F}(V) = \bigcap_{d=1}^{n-1}S_{d}
\end{equation}
is closed, so it's a projective variety.\par 

\begin{dang}
    \begin{proposition}
        We have $\dim (\Gr_d(V)) = d(n-d),\; n=\dim(V)$.
    \end{proposition}
\end{dang}
Proof: From equation~(\ref{equation: local description of a Grassmanian}), we see that at $U_{(1,\dots,d)}\subset \bbP(\bigwedge^dV)$, 
the Grassmanian $\Gr_d(V)$ is defined by $V(I)$, where
\begin{equation}
    I = (X_J-f_J(X):\forall J=(j_1,\dots,j_d)\in\mathcal{J}\setminus\{(1,\dots,d), J_{pq}: p,q\})
\end{equation}
Put any LEX order on $k[X_J:J\in\mathcal{J}\setminus\{(1,\dots,d)\}]$, with the property 
$X_J> X_{J_{pq}}$ whenever 
$J\notin \{J_{pq}: (p,q)\in\{1,\dots,d\}\times\{d+1,\dots,n\}\}$. Then, we see that $\sqrt{I}=I$ and 
\begin{equation}
    \frac{k[X_J:J\neq (1,\dots,d)]}{I}\cong k[X_{J_{pq}}: (p,q)\in\{1,\dots,d\}\times\{d+1,\dots,n\}]
\end{equation} 
which has transcendence degree $d(n-d)$.\qed

There is a simple way to represent an element $F\in \mathcal{F}(V)$. Fix a basis $\{e_1,\dots,e_n\}$ of $V$
and observe that we can find a basis of $V$ such that $F_i=\langle v_1,\dots, v_i\rangle$. Now, assume 
that $F\in U=U_{(1)}\times \cdots \times U_{(1,\dots,d-1)}\times U_{(1,\dots,d)}$. Then if we write 
$v_i = \sum_ja_{ij}e_j$, we see that $A=(a_{ij})$ has nonsingular principal minors $M_{ii},\;i=1\,\dots,n$. 
We call $A$ a representative matrix of $F$.\par  
Since $a_{11}\neq0$, we can assume $v_1=e_1+\sum_{j>1}a_{1j}e_j$. Now, observe that if we only do elementary operations 
of the form $R_j\to R_j + kR_i,\;(i<j)$ and $R_j\to kR_j,\;(k\neq0)$ (where $R_i=(a_{i1},\dots,a_{in})$ is the $i$th row of $A$) on $A$, we 
still have $A$ as a representative of the flag $F$. Using the hypothesis that the principal minors are nonsingular, we see that 
we can assume $a_{ii}=1,\forall i$, while $a_{ij}=0,\;(i>j)$. That is $v_i=e_i+\sum_{j>i}a_{ij}e_j$.

\begin{dang}
    \begin{proposition}
        We have $\dim\mathcal{F}(V)=\frac12 n(n-1)$.
    \end{proposition}
\end{dang}
Proof: The above discussion shows that $\mathcal{F}(V)\cap U$ is defined by the graph of a morphism  
$c: k^{\frac12 n(n-1)}\to k^N$ for some $N$. This proves the proposition.\qed

\begin{example}
Suppose that $F\in\mathcal{F}(k^3)$ has matrix representation 
\begin{equation}
    A = \begin{pmatrix}
        1&a&b\\ 0&1&c \\ 0&0&1
    \end{pmatrix}.
\end{equation}
Then, on the chart $U=U_{(1)}\times U_{(1,2)}\times U_{(1,2,3)}$ we have coordinates: 
\begin{equation}
    (\{(2)\mapsto a, (3)\mapsto b\}, \{(1,3)\mapsto c, (2,3)\mapsto ac - b\}, \{*\})
\end{equation} 
Then we see that $U\cap\mathcal{F}(k^3)$ is the graph of the morphism $k^3\to k: (a,b,c)\mapsto ac -b$.
\end{example}

\section{Lecture 9}
\subsection{Induced function spaces}
\subsection{General definition of product}

\section{Lecture 10}
In this section we let $X$ to be an abstract variety and $Y\subset X$ to be an irreducible subset 
(a point is a case of interest). Define 
\begin{equation}
    \mathcal{O}_{X,Y} = \varinjlim_{Y\cap U\neq\varnothing}\mathcal{O}_X(U)
\end{equation}
where the directed system is defined by $(\{\mathcal{O}_X(U)\}_{U\in \mathcal{U}}, \{\rho_{UV}\}_{U,V\in\mathcal{U}})$
where $\mathcal{U}$ is some system of generators for the topology of $X$. To avoid set theoretical issues we can 
take $\mathcal{U}$ to be the countable set of basic open sets of $X$. 

\begin{notation}
    We will denote the elements of $\mathcal{O}_{X,Y}$ by $\langle U,f\rangle$. So 
    $\langle U,f\rangle=\langle V,g\rangle$ if and only if there exists $W\subset U\cap V$ such that $f|_W=g|_W$. 
    In particular, since $Y$ is taken to be irreducible, we have that $W\cap Y \subset U\cap V\cap Y$ is dense, so 
    $f|_{U\cap V\cap Y}=g|_{U\cap V\cap Y}$ by continuity of $f$ and $g$.
\end{notation}

\begin{dang}
    \begin{proposition}
        The ring $\mathcal{O}_{X,Y}$ is a $k$-algebra and is a local ring.
    \end{proposition}
\end{dang}
Proof: The inclusion 
\begin{equation}
    k \hookrightarrow\mathcal{O}_{X,Y}:c\mapsto \langle X,c\rangle
\end{equation}
shows that $\callO_{X,Y}$ is a $k-$algebra. Now, define the set 
\begin{equation}
    \fr{m}_{X,Y} =\{\langle U,f\rangle: f|_{U\cap Y}=0\}.
\end{equation}
The condition $f|_{U\cap Y}=0$ is we defined on the class $\langle U,f\rangle$. In fact, if 
$V\subset U$ is open, then $f|_{V\cap Y}=0\Rightarrow f|_{U\cap Y}=0$, because $V\cap Y\subset U\cap Y$ is dense. Furthermore 
$\callO_{X,Y}$ is clearly seen to be an ideal. Finally, we observe that 
\begin{equation}
    \callO_{X,Y}-\fr{m}_{X,Y}\subset\callO_{X,Y}^*.
\end{equation}
In fact, if $\langle U,f\rangle$ is such that $f(U\cap Y)\neq\{0\}$, then the set 
$V=\{q\in U: f(q)\neq0\}$ is a nonempty open set of $U$. Therefore 
\begin{equation}
    \langle U,f\rangle = \langle V,f|_V\rangle = \langle V,\frac{1}{f}\bigg|_V\rangle ^{-1}\in\callO_{X,Y}^*.
\end{equation}
\qed 

\begin{dang}
    \begin{proposition}
        Let $X\subset k^m$ be an affine algebraic set and $Y\subset X$ be an irreducible subset. 
        Denote by $P=I(Y)/ I(X)\subset A = k[X]$, then $P$ is a prime ideal and we have an isomorphism
        \begin{equation}
          \psi:  A_P \rightarrow \callO_{X,Y}: \overline{F}/\overline{G}\mapsto \langle D_X(G), F/G\rangle.
        \end{equation}
    \end{proposition}
\end{dang}
Proof: First observe that $G\in I(Y)\iff G|_Y = 0\iff D_Y(G)=\varnothing$. Therefore 
$\overline{G}\notin P\iff Y\cap D(G)\neq\varnothing$.
Furthermore, we have $\overline{F}_1/\overline{G}_1 = \overline{F}_2/\overline{G}_2$ in $A_P$ 
if and only if there exists $G\notin P$ such that 
\begin{equation}
    \begin{split}
        &\overline{G}(\overline{F}_1\overline{G}_2 - \overline{F}_2\overline{G}_1) = 0 \in k[X]\iff\\
        &G(F_1G_2 - F_2G_1) |_X = 0 \in \callO_X(X)\iff\\
        &(F_1G_2 - F_2G_1)|_{D_X(G)} = 0\in\callO_X(D_X(G))\iff\\
        &\langle D_X(G), F_1/G_1\rangle = \langle D_X(G), F_2/G_2\rangle.
    \end{split}
\end{equation}
Then we see that $\psi$ is well defined and injective.
It remains to prove that it's surjective. To see this, let 
$\langle U, f\rangle \in\callO_{X,Y}$, and let $q\in U\cap Y$. Since $f:U\to k$ is regular, there exists 
$U^\prime\ni q, U^\prime\subset U$ and polynomials $F,G\in k[\underline{X}], G(q^\prime)\neq0\;\forall q^\prime\in U^\prime$ 
such that $f|_{U^\prime}=F/G |_{U^\prime}$. In particular, $q\in D_X(G)\cap Y$ and 
$\langle U,f\rangle = \langle D_X(G), F/G\rangle$.
To see that $P$ is prime, we observe that it's the image of the prime ideal $I(Y)$ under the surjective morphism 
$f:k[X_1,\dots,X_m]\to k[X]$ and $P\subset\ker(f)$.\qed

\begin{example}
    Let $X=k^1$ and $Y=\{0\}$, we have $A=k[t]$ and $P=I(Y)=(t)$. So we see that $A_P=k[t]_P =\{F/G: G\notin (t)\}$. 
    For the other side, we have 
    \begin{equation}
        \callO_{X,Y}=\{\langle D(G), F/G \rangle: G(0)\neq0\}
    \end{equation}
    because $D(G)\cap Y\neq\varnothing\iff G(0)\neq0$.
\end{example}

\section{Lecture 11}
\subsection{Birational geometry}

\section{Lecture 12}
\subsection{Blowup of affine varieties}
Consider the following data. 
\begin{itemize}
    \item $X\subset k^n$ an affine variety, $A=\callO_X(X)$ its coordinate ring;
    \item $I\subset A$ an ideal such that if $Y=V(I)$, then $U = X-Y\subset X$ is dense.
\end{itemize}
Denote $I=(f_1,\dots, f_m),f_i = F_i|_X$ and 
\begin{equation}
    h:k[X_1,\dots,X_n,Y_1,\dots,Y_m]\rightarrow A[T]:X_i\mapsto x_i = X_i|_X, Y_i\mapsto f_iT.
\end{equation}
Let $J=\text{Ker}(h)$, then $h$ is homogeneous in $\underline{Y}$. In fact, suppose that 
$F = F_0+F_2+\cdots\in J$, where $F_i$ the $i$th degree homogeneous component of $F$ viewed as a polynomial in 
$k[\underline{X}][\underline{Y}]$. Then 
\begin{equation}
    h(F_d) = F_d(x_1,\dots, x_n, f_1T,\dots, f_mT)=F_d(x_1,\dots,x_n,f_1,\dots,f_m)T^d\in A[T].
\end{equation}
So $h(F)=0\Rightarrow h(F_d)=0$ and $F_d\in J,\;\forall d\geq0$. 
\begin{remark}
    An important observation is that $F_iY_j - F_jY_i\in J,\forall i,j$ by construction.
\end{remark}
Now we can define 
\begin{equation}
    \widetilde{X} = V_{a|p}(J)\subset X\times\bbP^{m-1}\subset k^n\times\bbP^{m-1}.
\end{equation}
Observe that the first insclusion is a consequence of $p_1(\widetilde{X})\subset X$, where $p_1$ is the projection 
in $\underline{X}$. In fact, we have a injective map 
$p_1^\#:k[\underline{X}]\hookrightarrow k[\underline{X},\underline{Y}]:g\mapsto g\circ p_1$. Therefore 
$g\circ p_1\in J\iff g(x_1,\dots,x_n)=0\iff g|_X=0$, so $p_1^\# (I(X)) = J$ and $(a,b)\in\widetilde{X}$ implies that 
$g(a)=0,\;\forall g\in I(X)$, i.e., $a\in VI(X)=X$. 

\begin{remark}
By the discussion above, we see that 
$p_1^\#$ descends to a map $k[X]\to k[\underline{X},\underline{Y}]/J$.
\end{remark}

\begin{remark}\label{remark: U_i is affine}
    Denote by $U_i = \widetilde{X}\cap D(Y_i)$. Observe that $D(Y_i) = k^n\times D_+(Y_i)\subset k^n\times\bbP^m$ it's 
a open affine subset. So $U_i\subset D(Y_i)$ is closed, because it's the intersection of $D(Y_i)$ with a closed 
subset $\widetilde{X}\subset k^n\times\bbP^m$. As a closed subset of a affine subset, $U_i$ is an affine subset. 
In $U_i$ we have the relation $F_j = (\tfrac{Y_j}{Y_i})F_i$. Therefore 
\begin{equation}
    p_1^\#(I)\callO_{\widetilde{X}}(U_i) = (f_i\circ p_1|_{U_i}).
\end{equation}
\end{remark}


Now we define $E=p_1^{-1}(Y)=V_{a|p}(f_1\circ p_1,\dots, f_m\circ p_1)$ so we have the diagram.
\begin{center}
    \begin{tikzcd}
        E \arrow[r,hook]\arrow[d]&  \widetilde{X}\arrow[d,swap,"p_1"]\arrow[ddr,"p_2"] &\\ 
         Y\arrow[r,hook]&X\arrow[dr,dashed] & \\
         &&\bbP^m
    \end{tikzcd}
\end{center}
The dadhed arrow is the map $p\mapsto (f_1(p):\cdots:f_m(p))$, it's not defined only on $Y$ and since 
$X-Y$ is dense in $X$ by hypothesis, it's a rational map.
\begin{definition}
    The Rees algebra of $I\subset A$ is defined as the image the map $h$ defined below. That is, we let 
    \begin{equation}
        \text{Rees}_I(A) = A\oplus IT\oplus I^2T^2\oplus\cdots\oplus I^dT^d\oplus\cdots \subset A[T].
    \end{equation}
\end{definition}
\begin{remark}
    Remember that $I =\bigcap_{P\supset I,\text{ prime}}P$. We say that $A$ is reduced if $\sqrt{(0)}=(0)$, i.e, 
    $x^n=0\Rightarrow x=0$. In particular, this means that $(0)=\bigcap_{P\text{ prime}}P$. Now, assume that 
    $A$ is a reduced ring (it's the case of $A=\callO_X(X)$), we claim that $A[T]$ is reduced. In fact, 
    $f^n=0\in A[T]$ implies that $f^n=0\in (\tfrac{A}{P})[T]$ for all $P$ prime. Since $A/P$ is a domain, we see that 
    $f=0\in (\tfrac{A}{P})[T]$ for any prime $P$. So if $f=a_0+a_1T+\cdots$, we get that $a_i\in P$ for all prime $P$ for all
    $i\geq0$. Therefore $a_i=0,;\forall i$ because $A$ is reduced.
\end{remark}

By the previous remark, we see that $\text{Rees}_I(A)$ is reduced and since $k[\underline{X},\underline{Y}]/J\cong\text{Rees}_I(A)$, 
we get that $J$ is radical ($J=\sqrt{J}$). 
\begin{dang}
    \begin{lemma}
        If $Z\subset k^m$ is affine and $f\in \callO_Z(Z)$ is such that $f$ is not a divisor of zero in 
        $\callO_Z(Z)$, then $Z-V(f)\subset Z$ is dense.
    \end{lemma}
\end{dang}
Proof: If it's not dense, we get that $D_Z(g)\subset V(f)$ for some $g$. So $fg|_Z=0$ and $f|_Z$ is a divisor of zero.\qed 
\begin{dang}
\begin{proposition}
    $f_i\circ p_1|_{U_i}$ is not a divisor of zero.
\end{proposition}
\end{dang}
Proof: Let $b\in\callO_{\widetilde{X}}(U_i)\setminus\{0\}$ such that $b\cdot f_i\circ p_1|_{U_i}=0$ (by contradiction). 
Since $U_i$ is affine (see remark~(\ref{remark: U_i is affine})), we have $b=B/Y_i|_{U_i}$ with $B\in k[\underline{X},\underline{Y}]$ homogeneous of degree 
$d$ in $Y$. Since $b\neq0$, we have $B\notin J$. Now 
\begin{equation}
    \frac{B}{Y_i^d}\cdot F_i\bigg|_{\widetilde{X}\cap D(Y_i)}=0\Rightarrow BF_iY_i\big|_{\widetilde{X}}=0
\end{equation}
which implies that $BF_iY_i \in \sqrt{J}=J$. Therefore
\begin{equation}
    \begin{split}
        &B(x_1,\dots,x_n, f_1T,\dots, f_mT)f_i^2 T = 0\\
        &\Rightarrow B(x_1,\dots,x_n,f_1,\dots,f_m)f_i^2T = 0\\
        &\Rightarrow B(x_1,\dots, x_n, f_1,\dots, f_m)|_{D_X(f_i)} = 0\\
        &\Rightarrow B(x_1,\dots, x_n, f_1,\dots, f_m)|_{X-Y}=0\;(\text{Why ??})\\
        &\Rightarrow B(x_1,\dots,x_n,f_1,\dots, f_m) = 0\\
        &\Rightarrow B(X_1,\dots, X_n, Y_1,\dots, Y_m)\in J
    \end{split}
\end{equation}
which is a contradiction.\qed
\begin{dang}
    \begin{proposition}
        $\widetilde{X}-E\subset \widetilde{X}$ is dense.
    \end{proposition}
\end{dang}
Proof: By the last proposition, we see that $U_i - E\cap U_i = U_i - V(f_i\circ p_1)$ is dense in $U_i$ and since 
$E\subset \widetilde{X}=\bigcup_i U_i$, we get that $\widetilde{X}-E$ is dense in $\widetilde{X}$.\qed

\underline{Summary}: We constructed the following data. 
\begin{center}
    \begin{tikzcd}
        E =b^{-1}(Y)\arrow[r,hook]\arrow[d] & \widetilde{X}\arrow[d,"b"] & \widetilde{X} -E\arrow[l,hook,swap,"\text{dense}"]\arrow[d,"s"]\\         
        Y\arrow[r,hook]&   X  & X- Y\arrow[l,hook,"\text{dense}"]
    \end{tikzcd}
\end{center}
Furthermore, there exists a covering $\widetilde{X}=\bigcup_i U_i$ of open affines, sections 
$b_i\in\callO_X(U_i),\;b_i\nmid 0$ such that $I\callO_X(U_i)=(b_i)$.
\subsection{universal property}

\subsection{Examples}
\subsection{Generalization}